% !TEX program = xelatex
\documentclass[11pt]{article}
\usepackage{fontspec}
\usepackage{xeCJK}
\usepackage[margin=1in]{geometry}
\usepackage{amsmath}
\usepackage{amssymb}
\usepackage{amsthm}
\usepackage{graphicx}
\usepackage{tikz-cd}
\usepackage{multicol}
\usepackage{titling}
\usepackage{setspace}
\usepackage{float}
\usepackage{subcaption}
\usepackage{algorithm}
\usepackage{algpseudocode}
\newtheorem{theorem}{Theorem}
\newtheorem{lemma}{Lemma}

% Bigger & bolder title
\pretitle{\begin{center}\LARGE\bfseries}
\posttitle{\par\end{center}\vskip 0.5em}
\predate{\begin{center}}
\postdate{\end{center}}
\setlength{\droptitle}{-6em}

% Paragraph style (old)
% \setlength{\parindent}{0pt}
% \setlength{\parskip}{0.4\baselineskip}
% Paragraph style (new)
% \setlength{\parindent}{2em}
% \setlength{\parskip}{0pt}
% \setstretch{0.96}


\begin{document}

\title{Notes on Coding Theory}
\author{}
\date{}
\maketitle

\clearpage
\section{Low-Density Parity-Check Codes}

Low-density parity-check (LDPC) codes are a class of linear block codes with implementable decoders, which provide near-capacity performance on a large set of data-transmission and data-storage channels.

\subsection{Representation of LDPC Codes}

A low-density parity-check code is a linear block code given by the null space of an $m \times n$ parity-check matrix \textbf{H} that has a low density of 1s. The density need only be sufficiently low to permit effective iterative decoding.

\[
    \mathbf{H} \in \{0,1\}^{m \times n}, \qquad \mathbf{H}\mathbf{t} = \mathbf{0} \quad \text{for all codewords } \mathbf{t}
\]

LDPC codes can be classified as:
\begin{itemize}
  \item A regular LDPC code is a linear block code whose parity-check matrix \textbf{H} has fixed column weight $g$ and row weight $r$, where $r/g = n/m$ and $g \ll m$. In this case, the code rate $R$ is bounded by $R \geq 1 - \frac{m}{n} = 1 - \frac{g}{r}$ with equality when \textbf{H} is full rank.
  \item An irregular LDPC code is a linear block code whose parity-check matrix \textbf{H} has variable column weights and row weights, being determined by code-design procedures.
\end{itemize}

\noindent Additionally, almost all LDPC code constructions impose the structural property on \textbf{H}: no two rows or two columns have more than one position in common that contains a nonzero element.

Graphically, \textbf{Tanner graphs} are used to represent a LDPC code from its parity-check matrix \textbf{H}. Tanner graph is a bipartite graph, that is, a graph whose nodes may be separated into two types, with edges connecting only nodes of different types. The two types of nodes in a Tanner graph are the variable nodes (or code-bit nodes) and the check nodes (or constraint nodes), which are denoted by VNs and CNs, respectively.

In a Tanner graph, CN i is connected to VN j whenever element $h_{ij}$ in \textbf{H} is a $1$. Therefore, there are m CNs in a Tanner graph, one for each check equation, and n VNs, one for each code bit; the m rows of \textbf{H} specify the m CN connections, and the n columns of \textbf{H} specify the n VN connections. This is the basis of \textbf{iterative decoding}: each of the nodes acts as a locally operating processor and each edge acts as a bus that conveys information from a given node to each of its neighbors.

\begin{figure}[H]
  \centering
  \includegraphics[width=0.5\linewidth]{Tannerdemo.png}
  \caption{Tanner graph of a (10,5) linear block code with $r = 4$, $g = 2$ (i.e. the degree of each VN is 2 and the degree of each CN is 4)}
  \label{fig:tanner-demo}
\end{figure}

The parity-check matrix for the Tanner graph above is 
\[
\mathbf{H} =
\begin{bmatrix}
1 & 1 & 1 & 1 & 0 & 0 & 0 & 0 & 0 & 0 \\
1 & 0 & 0 & 0 & 1 & 1 & 1 & 0 & 0 & 0 \\
0 & 1 & 0 & 0 & 1 & 0 & 0 & 1 & 1 & 0 \\
0 & 0 & 1 & 0 & 0 & 1 & 0 & 1 & 0 & 1 \\
0 & 0 & 0 & 1 & 0 & 0 & 1 & 0 & 1 & 1
\end{bmatrix}.
\]
As follows from $\mathbf{v}\mathbf{H}^{T} = \mathbf{0}$, that the bit values connected to the same check node must sum to zero (mod 2). Also, since rank(\textbf{H}) $= 4 \neq 5$, \textbf{H} is not full rank, so the code rate is $R = 1- 4/10 = 3/5$. As shown in \ref{fig:tanner-demo}, the six thickened edges form a \textbf{closed path}, which is called a \textbf{cycle}. The length of a cycle is equal to the number of edges which form the cycle: a cycle of length l is often called on l-cycle, and the minimum cycle length in a given bipartite graph is called the \textbf{graph's girth}.

It is possible to more closely approach capacity limits with irregular LDPC codes than with regular LDPC codes. For irregular LDPC codes, the parameters g and r vary with the columns and rows, so it is usual to specify the VN and CN \textbf{degree-distribution polynomials}. In the polynomial,
\[
\lambda(X) = \sum_{d=1}^{d_v} \lambda_d X^{d-1},
\]
$\lambda_d$ denotes the fraction of all edges connected to degree-$d$ VNs and $d_v$ denotes the maximum VN degree. Similarly, in the polynomial
\[
\rho(X) = \sum_{d=1}^{d_c} \rho_d X^{d-1},
\]
$\rho_d$ denotes the fraction of all edges connected to degree-$d$ CNs and $d_c$ denotes the maximum CN degree. Note that, for the regular code above, for which $g = 2$ and $r = 4$, we have $\lambda(X) = X$ and $\rho(X) = X^3$. It is also possible to view this from a node perpective, $L(X) = \sum_{d=1}^{d_v} L_d X^{d}$ and $R(X) = \sum_{d=1}^{d_c} R_d X^{d}$, where $L_d = \frac{\lambda_d/d}{\int_0^1 \lambda(X)\,dX}$ and $R_d = \frac{\rho_d/d}{\int_0^1 \rho(X)\,dX}$.

The average degrees are $\bar{d}_v = L'(1) = \left(\int_0^1 \lambda(x)\,dx\right)^{-1}$ and $\bar{d}_c = R'(1) = \left(\int_0^1 \rho(x)\,dx\right)^{-1}$. The design rate of an LDPC code is $R \geq 1 - \frac{\bar{d}_v}{\bar{d}_c} = 1 - \frac{L'(1)}{R'(1)} = 1 - \frac{\int_0^1 \rho(x)\,dx}{\int_0^1 \lambda(x)\,dx}$. 

\clearpage
\subsection{Message Passing and the Turbo Principle}

\textbf{Message-passing decoding} refers to a collection of low-complexity decoders working in a distributed fashion to decode a received codeword in a concatenated coding scheme.

\begin{figure}[H]
  \centering
  \includegraphics[width=0.5\linewidth]{turboprinciple.png}
  \caption{A depiction of the trubo principle for concatenated coding schemes.(a) Extrinsic information $I_{X \to Y}$ from decoder X to decoder Y. (b) Extrinsic information $I_{Y \to X}$ from decoder Y to decoder X. (c) A compact symmetric combination of (a) and (b).}
  \label{fig:turbo}
\end{figure}

The notion of \textbf{extrinsic-information passing} is called the turbo principle in the context of the iterative decoding of concatenated codes in communication channels. \textbf{Total information}

\[
I^{\text{total}}_X = \sum_{Z \in N(X)} I_{Z \to X}  + I_X
\]
and the \textbf{extrinsic information}
\[
I_{X \to Y} = \sum_{Z \in N(X) - \{Y\}} I_{Z \to X}  + I_X
\]
where $I_X$ denotes the \textbf{intrinisic information} received from the channel. Essentially, Y receives the sum total of messages that
X has available minus the information that Y already has, which means X does not pass to a neighboring decoder any information that the neighboring decoder already has.

\clearpage
\subsection{Sum-Prudct Algorithm}
The sum-product algorithm (SPA) or belief propagation algorithm (BPA) for general memoryless binary-input channels applies the turbo principle to correct the errors present in the received code bits.

In the context LDPC code, it can be considered as \textit{a generalized concatenation of many Single-Parity-Check code (CNs) and Repetition code (VNs)}, and they work cooperatively in a distributed fashion to determine the correct code bit values.

\begin{figure}[H]
  \centering
  \includegraphics[width=0.5\linewidth]{LDPCgraph.png}
  \caption{Graphical representation of an LDPC code as a concatenation of SPC and REP codes. $\pi$ represents an interleaver.}
  \label{fig:LDPC-demo}
\end{figure}

The optimality criterion underlying the development of the SPA decoder is symbol-wise maximum a posteriori (MAP). we are interested in computing the a posteriori probability (APP) that a specific bit in the
transmitted codeword $\mathbf{v} = [v_0 \; v_1 \; \cdots \; v_{n-1}]$ equals to 1, given the received codeword $\mathbf{y} = [y_0 \; y_1 \; \cdots \; y_{n-1}]$. For a general code bit $v_j$, the APP is 
\[
\Pr(v_j = 1 \mid \mathbf{y})
\]
the APP ratio (or likelihood ratio under the uniform priors assumption),
\[
l(v_j \mid \mathbf{y}) \triangleq \frac{\Pr(v_j = 0 \mid \mathbf{y})}{\Pr(v_j = 1 \mid \mathbf{y})}
\]
or log-likelihood ratio (LLR) can be used for numerical stability,
\[
L(v_j \mid \mathbf{y}) \triangleq 
\log\left(\frac{\Pr(v_j = 0 \mid \mathbf{y})}{\Pr(v_j = 1 \mid \mathbf{y})}\right)
\]

The VN and CN decoders work cooperatively and iteratively to estimate $L(v_j \mid \mathbf{y})$ for $j = 0, 1, . . ., n-1$. Throughout our development, the \textbf{flooding schedule} is employed. According to this schedule, all VNs process their inputs and pass extrinsic information up to their neighboring check nodes; the check nodes then process their inputs and pass extrinsic information down to their neighboring variable nodes; and the procedure repeats, starting with the variable nodes. After a preset maximum number of iterations of this VN/CN decoding round, or after some stopping criterion has been met, the decoder computes (estimates) the LLRs $L(v_j \mid \mathbf{y})$ from which decisions on the bits $v_j$ are made. When the cycles are large, the estimates will be very accurate and the decoder will have near-optimal (MAP) performance. 

SPA assumes the \textit{LLR quantities received at each node from its neighbors are independent}. Clearly this breaks down when the number of iterations exceeds half of the Tanner graph's girth.

\subsubsection{VN: Repetition Code MAP Decoder and APP Processor}

\begin{figure}[H]
  \centering
  \includegraphics[width=0.5\linewidth]{VNpassing.png}
  \caption{A VN decoder (REP decoder). VN $j$ receives LLR information from the channel and from all of its neighbors, excluding message $L_{i \to j}$ from CN $i$, from which VN $j$ composes message $L_{j \to i}$ that is sent to CN $i$}
  \label{fig:VN-demo}
\end{figure}

Consider a REP code in which the binary code symbol $c \in \{0,1\}$ is transmitted over a memoryless channel $d$ times so that the $d$-vector $\mathbf{r}$ is received. By definition, the MAP decoder computes the log-APP ratio
\[
L(c \mid \mathbf{r}) = \log\left(\frac{\Pr(c = 0 \mid \mathbf{r})}{\Pr(c = 1 \mid \mathbf{r})}\right) = \log\left(\frac{\Pr(\mathbf{r} \mid c = 0)}{\Pr(\mathbf{r} \mid c = 1)}\right)
\]
under an equally likely assumption for the value of $c$ (ratio of a posterior = ratio of likelihood)

This simplifies as
\[
\begin{aligned}
L(c \mid \mathbf{r})
&= \log\left(\frac{\prod_{l=0}^{d-1} \Pr(r_l \mid c = 0)}{\prod_{l=0}^{d-1} \Pr(r_l \mid c = 1)}\right) \\
&= \sum_{l=0}^{d-1} \log\left(\frac{\Pr(r_l \mid c = 0)}{\Pr(r_l \mid c = 1)}\right) = \sum_{l=0}^{d-1} L(r_l \mid x)
\end{aligned}
\]
where $L(r_l \mid x)$ is obviously defined. Thus, the MAP receiver for a REP code computes the LLRs for each channel output $r_l$ and adds them. \textit{The MAP decision is $\hat{c} = 0$ if $L(c \mid \mathbf{r}) \geq 0$ and $\hat{c} = 1$ otherwise.}

In the context of LDPC decoding, the above expression is adapted to compute the extrinsic information to be sent from VN $j$ to CN $i$ of Figure~\ref{fig:VN-demo},
\[
\boxed{
L_{j \to i} = L_j + \sum_{i' \in N(j) - \{i\}} L_{i' \to j}
}
\]
The quantity $L_j$ in this expression is the LLR value computed from the channel sample $y_j$,
\[
L_j = L(c_j \mid y_j) = \log\left(\frac{\Pr(c_j = 0 \mid y_j)}{\Pr(c_j = 1 \mid y_j)}\right).
\]
At the last iteration, VN $j$ produces a decision based on
\[
L_j^{\text{total}} = L_j + \sum_{i \in N(j)} L_{i \to j}
\]

\subsubsection{CN: Single-Parity-Check Code MAP Decoder and APP Processor}

\begin{figure}[H]
  \centering
  \includegraphics[width=0.5\linewidth]{CNpassing.png}
  \caption{A CN decoder (SPC decoder). CN $i$ receives LLR information from all of its neighbors, excluding message $L_{j \to i}$ from VN $j$, from which CN $i$ composes message $L_{i \to j}$ that is sent to VN $j$}
  \label{fig:CN-demo}
\end{figure}

\begin{lemma}
\label{lem:gallager-even-odd}
  Consider a vector of $d$ independent binary random variables $\mathbf{a} = [a_0, a_1, \ldots, a_{d-1}]$ in which $\Pr(a_l = 1) = p_1^{(l)}$ and $\Pr(a_l = 0) = p_0^{(l)}$. Then the probability that $\mathbf{a}$ contains an even number of 1s is
\[
\frac{1}{2} + \frac{1}{2}\prod_{l=0}^{d-1}\left(1 - 2p_1^{(l)}\right),
\]
and the probability that $\mathbf{a}$ contains an odd number of 1s is
\[
\frac{1}{2} - \frac{1}{2}\prod_{l=0}^{d-1}\left(1 - 2p_1^{(l)}\right).
\]
\end{lemma}

Consider the transmission of a length-$d$ SPC codeword $\mathbf{c}$ over a memoryless channel whose output is $\mathbf{r}$. The bits in the codeword $\mathbf{c}$ have a single constraint: \textit{there must be an even number of 1s in $\mathbf{c}$}. Without loss of generality, we focus on bit $c_0$, for which the MAP decision rule is
\[
\hat{c}_0 = \arg\max_{b \in \{0,1\}} \Pr(c_0 = b \mid \mathbf{r}, \mathrm{SPC}),
\]
where the conditioning on SPC is a reminder that there is an SPC constraint imposed on $\mathbf{c}$. Consider now
\[
\begin{aligned}
\Pr(c_0 = 0 \mid \mathbf{r}, \mathrm{SPC})
&= \Pr(c_1, c_2, \ldots, c_{d-1} \text{ has an even number of 1s} \mid \mathbf{r}) \\
&= \frac{1}{2} + \frac{1}{2}\prod_{l=1}^{d-1}\left(1 - 2\Pr(c_l = 1 \mid r_l)\right),
\end{aligned}
\]
where the second line follows from the \ref{lem:gallager-even-odd}. Rearranging gives
\[
1 - 2\Pr(c_0 = 1 \mid \mathbf{r}, \mathrm{SPC}) = \prod_{l=1}^{d-1}\left(1 - 2\Pr(c_l = 1 \mid r_l)\right),
\]
where $p(c_0 = 0 \mid \mathbf{r}, \mathrm{SPC}) = 1 - p(c_0 = 1 \mid \mathbf{r}, \mathrm{SPC})$. 

We can change this to an \textbf{LLR representation} using the easily proven relation for a generic binary random variable with probabilities $p_1$ and $p_0$,
\[
1 - 2p_1 = \tanh\left(\frac{1}{2}\log\left(\frac{p_0}{p_1}\right)\right) = \tanh\left(\frac{1}{2}\mathrm{LLR}\right).
\]
where here $\mathrm{LLR} = \log(p_0/p_1)$. Applying this relation gives
\[
\tanh\left(\frac{1}{2}L(c_0 \mid \mathbf{r}, \mathrm{SPC})\right)
= \prod_{l=1}^{d-1}\tanh\left(\frac{1}{2}L(c_l \mid r_l)\right)
\]
or
\[
L(c_0 \mid \mathbf{r}, \mathrm{SPC})
= 2\tanh^{-1}\left(\prod_{l=1}^{d-1}\tanh\left(\frac{1}{2}L(c_l \mid r_l)\right)\right).
\]
Thus, \textit{the MAP decoder for bit $c_0$ in a length-$d$ SPC code makes the decision $\hat{c}_0 = 0$ if $L(c_0 \mid \mathbf{r}, \mathrm{SPC}) \geq 0$ and $\hat{c}_0 = 1$ otherwise.}

In the context of LDPC decoding, when the CNs function as APP processors instead of MAP decoders, CN $i$ computes the extrinsic information
\[
\boxed{
L_{i \to j} = 2\tanh^{-1}\left(\prod_{j' \in N(i)-\{j\}} \tanh\left(\frac{1}{2}L_{j' \to i}\right)\right)
}
\]
and transmits it to VN $j$. 

\subsubsection{The Gallager SPA Decoder}

\begin{algorithm}[H]
\caption{The Gallager Sum-Product Algorithm}
\begin{algorithmic}[1]
\State \textbf{Initialization:} For all $j$, initialize $L_j$ according to the appropriate channel model. Then, for all $i,j$ for which $h_{ij} = 1$, set $L_{j \to i} = L_j$.
\State \textbf{CN update:} Compute outgoing CN messages $L_{i \to j}$ for each CN using
  \[
  L_{i \to j} = 2\tanh^{-1}\left(\prod_{j' \in N(i)-\{j\}} \tanh\left(\frac{1}{2}L_{j' \to i}\right)\right),
  \]
  and then transmit them to the VNs.
\State \textbf{VN update:} Compute outgoing VN messages $L_{j \to i}$ for each VN using
  \[
  L_{j \to i} = L_j + \sum_{i' \in N(j)-\{i\}} L_{i' \to j},
  \]
  and then transmit them to the CNs.
\State \textbf{LLR total:} For $j = 0,1,\ldots,n-1$, compute
  \[
  L_j^{\text{total}} = L_j + \sum_{i \in N(j)} L_{i \to j}.
  \]
\State \textbf{Stopping criteria:} For $j = 0,1,\ldots,n-1$, set
  \[
  \hat{v}_j =
  \begin{cases}
  1, & \text{if } L_j^{\text{total}} < 0, \\
  0, & \text{otherwise,}
  \end{cases}
  \]
  to obtain $\hat{\mathbf{v}}$. If $\hat{\mathbf{v}}\mathbf{H}^{T} = \mathbf{0}$ or the number of iterations equals the maximum limit, stop; otherwise, go to Step 2.
\end{algorithmic}
\end{algorithm}

The decoder is initialized by setting all VN messages $L_{j \to i}$ equal to
\[
L_j = L(v_j \mid y_j) = \log\left(\frac{\Pr(v_j = 0 \mid y_j)}{\Pr(v_j = 1 \mid y_j)}\right),
\]
for all $j,i$ for which $h_{ij} = 1$. Here, $y_j$ represents the channel value that was actually received, that is, it is not a variable here. We consider the following special cases.

\textbf{BEC.} In this case, $y_j \in \{0,1,e\}$ and we define $p = \Pr(y_j = e \mid v_j = b)$ to be the erasure probability, where $b \in \{0,1\}$. Then it is easy to see that
\[
\Pr(v_j = b \mid y_j) =
\begin{cases}
1, & \text{when } y_j = b, \\
0, & \text{when } y_j = b^c, \\
1/2, & \text{when } y_j = e,
\end{cases}
\]
where $b^c$ represents the complement of $b$. From this, it follows that
\[
L(v_j \mid y_j) =
\begin{cases}
+\infty, & y_j = 0, \\
-\infty, & y_j = 1, \\
0, & y_j = e.
\end{cases}
\]

\textbf{BSC.} In this case, $y_j \in \{0,1\}$ and we define $\varepsilon = \Pr(y_j = b^c \mid v_j = b)$ to be the error probability. Then it is obvious that
\[
\Pr(v_j = b \mid y_j) =
\begin{cases}
1-\varepsilon, & \text{when } y_j = b, \\
\varepsilon, & \text{when } y_j = b^c.
\end{cases}
\]
From this, we have
\[
L(v_j \mid y_j) = (-1)^{y_j}\log\left(\frac{1-\varepsilon}{\varepsilon}\right).
\]
Note that knowledge of $\varepsilon$ is necessary.

\textbf{BI-AWGNC.} We first let $x_j = (-1)^{v_j}$ be the $j$th transmitted binary value; note that $x_j = +1(-1)$ when $v_j = 0(1)$. We shall use $x_j$ and $v_j$ interchangeably hereafter. The $j$th received sample is $y_j = x_j + n_j$, where the $n_j$ are independent and normally distributed as $\mathcal{N}(0,\sigma^2)$. Then it is easy to show that
\[
\Pr(x_j = x \mid y_j) = \left[1 + \exp\left(-2y_jx/\sigma^2\right)\right]^{-1},
\]
where $x \in \{\pm 1\}$ and, from this, that
\[
L(v_j \mid y_j) = 2y_j/\sigma^2.
\]
In practice, an estimate of $\sigma^2$ is necessary.

\textbf{Rayleigh.} Assuming an independent Rayleigh fading channel and the decoder has perfect channel state information. The model is similar to that of the AWGNC: $y_j= \alpha_j x_j + n_j$, where $\{ \alpha_j \}$ are independent Rayleigh random variables with unity variance. The channel transition probability in this case is
\[
\Pr(x_j = x \mid y_j) = \left[1 + \exp\left(-2\alpha_j y_j x/\sigma^2\right)\right]^{-1}.
\]
Then the variable nodes are initialized by
\[
L(v_j \mid y_j) = 2\alpha_j y_j/\sigma^2.
\]
In practice, estimates of $\alpha_j$ and $\sigma^2$ are necessary.

\clearpage
\subsection{Generalized LDPC Codes}
In a standard LDPC code, the two types of nodes in the Tanner graph perform 
simple operations. \textbf{Variable nodes (VN)} act as \textbf{repetition (REP) 
codes}, repeating the same bit value across all connected edges --- a degree-$d_v$ 
variable node is equivalent to a rate-$1/d_v$ repetition code. \textbf{Check nodes 
(CN)} act as \textbf{single parity-check (SPC) codes}, enforcing that the XOR of 
all connected bits equals zero. A standard LDPC code can therefore be viewed as a 
graph-based concatenation of REP and SPC component codes.

A \textbf{Generalized LDPC (GLDPC)} code replaces these simple component codes with 
stronger linear block codes. Check nodes are replaced by arbitrary linear block codes 
(e.g.\ Hamming, BCH, or Reed--Solomon codes) that enforce richer constraints than a 
single parity check, while variable nodes may similarly be replaced by stronger codes 
beyond simple repetition. The belief propagation (BP) decoder naturally extends to 
GLDPC codes: the simple REP and SPC message-passing rules at each node are replaced 
by the \textbf{a posteriori probability (APP) processor} for the corresponding 
component code, while the overall graph structure and message-passing framework remain unchanged.

The primary advantage of GLDPC codes is improved \textbf{minimum distance} for finite 
block lengths, which directly reduces the error floor, along with faster convergence 
due to stronger per-iteration information exchange. The cost is increased decoding 
complexity per node, since each check node now runs a full sub-decoder rather than a 
simple XOR operation.

\end{document}